\begin{exercisechunk}
\item \textbf{Logits, sigmoid probabilities, and logistic loss.}
\label[exercise]{ex:l10-logistic-log-loss}
\exerciseprofile{2}{2}{2}{\exproof\quad\excalculation\quad\exinterpretation}{%
Derive binary logistic loss from a conditional probability model and
distinguish the probabilistic surrogate from classification error.}
Define the sigmoid function by
\[
 \sigma(s)=\frac1{1+e^{-s}},
 \qquad s\in\R.
\]
For a score $s_w(x)=\langle w,x\rangle$ and a label
$y\in\{-1,1\}$, define
\[
 q_w(y\mid x)=\sigma\bigl(ys_w(x)\bigr).
\]
Lecture~1 introduced a binary logit as log-odds, and Lecture~3 introduced
sigmoid probabilities and log-loss fitting. This exercise reconnects those
ideas for a linear score.
\begin{enumerate}[(a),beginpenalty=10000]
\item Prove that $\sigma(s)+\sigma(-s)=1$. Deduce that
$q_w(\cdot\mid x)$ is a probability mass function on $\{-1,1\}$ for every
$x$.

\item Prove the log-odds identity
\[
 \log\frac{q_w(1\mid x)}{q_w(-1\mid x)}=s_w(x),
\]
and derive
\[
 -\log q_w(y\mid x)
 =\log\bigl(1+e^{-ys_w(x)}\bigr).
\]
Use these identities to explain each part of the terminology: why $s_w(x)$ is
a \emph{logit}, why the use of $\sigma$ makes the model \emph{logistic}, why
modeling $\P(Y=1\mid X=x)$ as a function of $x$ is called \emph{regression},
and why the negative log-likelihood is called \emph{logistic loss}.

\item Fix an input $x$ and suppose
$\pi=\P(Y=1\mid X=x)\in(0,1)$. If a predictor assigns probability
$q\in(0,1)$ to label $1$, its conditional expected log loss is
\[
 C_\pi(q)=-\pi\log q-(1-\pi)\log(1-q).
\]
Prove that $C_\pi$ is uniquely minimized at $q=\pi$. State precisely what the
identity $q=\pi$ implies when the logistic model can represent the true
conditional probabilities.

\item Show that $q_w(1\mid x)\ge1/2$ if and only if $s_w(x)\ge0$.
Compare the logistic training loss with the deployment loss
$\I\{\operatorname{sign}(s_w(x))\ne y\}$. Why is logistic loss called a
surrogate loss for classification error?
\end{enumerate}
\begin{hints}
For (c), differentiate $C_\pi(q)$ and determine the sign of its derivative
on either side of $q=\pi$. For (d), the two losses use the same signed score,
but one records its magnitude and the other records only whether its sign is
correct.
\end{hints}
\end{exercisechunk}
